\label{chap.2}

This chapter positions the work. Section~\ref{sec:bg-webtesting} covers automated web
testing and its oracles, Section~\ref{sec:bg-rltesting} RL for web and GUI testing,
Section~\ref{sec:bg-llm} LLM-driven testing and LLM-derived rewards,
Section~\ref{sec:bg-drl} the deep RL and exploration methods the testers are built from,
and Section~\ref{sec:bg-eval} how such tools are evaluated. Section~\ref{sec:bg-gap} states
the gap this thesis addresses. Statements about what a paper did or found are taken from its
abstract or full text; paragraphs headed \emph{For this work} are our reasoning.

\section{Automated web testing and test oracles}
\label{sec:bg-webtesting}

An automated web tester drives a browser through an application and decides, at each page,
what to do next. Crawljax established the model-based approach for
JavaScript-heavy applications: it infers a state-flow graph of user-interface states by
firing events and comparing the resulting DOM. Model-based and search-based testers explore
systematically, while random ("monkey") testers are cheap and, in GUI testing, a notoriously
strong baseline.

Exploring is only half of testing. The other half is deciding whether what happened was
wrong: the oracle problem. Barr et al.~\cite{barr2015oracle} classify oracles as specified
(a specification says what is correct), derived (from artefacts such as earlier versions),
implicit (generic failures such as crashes and exceptions) and absent (left to a human). RL
web testers use implicit oracles, namely errors in the browser console~\cite{zheng2021webexplor, gu2025webrled}. These are cheap and general, but blind to wrong content that renders without
error.

A \emph{seeded-fault oracle} is a specified oracle for a known set of faults. It provides
ground truth, which evaluation methodology in fuzzing prefers: Klees
et al.~\cite{klees2018evaluating} recommend measuring against known bugs rather than
heuristic proxies, and Magma inserts real bugs into real software with
instrumentation so that bug-centric metrics do not depend on the tool. The cost is validity.
Just et al. studied 357 real faults in five Java programs and found
mutant detection significantly correlated with real-fault detection, but 17\% of real
faults were not coupled to any mutant. Seeded faults are a supported but incomplete proxy.

Once a failure is found, delta debugging~\cite{zeller2002simplifying} minimises the failing
input while preserving the failure. Its original case study reduced 95 user actions that
crashed a browser to 3; it is used in the same way here, to shorten a reproduction trace.

\section{Reinforcement learning for web and GUI testing}
\label{sec:bg-rltesting}

\paragraph{Web.} WebExplor~\cite{zheng2021webexplor} is the reference point. It uses
curiosity-driven RL to generate action sequences, with tabular Q-learning and a curiosity
reward computed from a visit-count table over transitions, and builds an automaton online
that sends the agent back to under-explored states. Training and testing are intertwined:
the policy is learned on the application under test. Its failures are console errors, and
the authors note that this oracle may miss UI bugs. WebQT changes the
reward to encourage interaction like a human tester's and treats pages with the same
function as one state, reporting higher code coverage on seven open-source applications.
QExplore combines tabular Q-learning with a guided search and text
input generation.

WebRLED~\cite{gu2025webrled} combines a deep Q-network with a grid-based action-value
scheme, a learned action discriminator, and a curiosity reward that considers novelty both
within an episode and over the global history. Liu et al. train a DQN
or PPO agent to guide a crawler towards higher branch coverage. MARG runs
several Q-learning agents that share Q-table data; its successor
WebCQ moves to deep multi-agent RL and builds a vector for each
candidate action, concatenated with the state vector, because "the available UI events vary
across states". WebCQ is the closest architectural neighbour of this thesis's first tester.

A 2026 empirical study~\cite{liu2026webgui} compares model-based, RL-based and LLM-based
web testers under six state abstractions. No single strategy excels on every dimension;
RL-based approaches most often achieve the highest coverage; compact state abstractions suit
RL; and code coverage is only weakly correlated with the ability to reveal failures.

\paragraph{Mobile.} Q-testing uses a curiosity reward from a neural
state-comparison module. ARES reports that deep RL achieves higher
coverage and fault revelation than baselines including Q-testing.
DinoDroid~\cite{zhao2024dinodroid} is the clearest precedent for learning across
applications: a deep Q-network learns a behaviour model from a set of existing apps and is
used on new ones, evaluated by two-fold cross-validation over 64 apps (32 for training in
each fold). Humanoid~\cite{li2019humanoid} learns from human interaction traces and reports
higher coverage than six other generators, but its authors add that the learned model "was
not the main reason" for the improvement.

\paragraph{For this work.} The RL agent design is not new in itself. Count-based curiosity
is the default reward in this line of work, per-action scoring over a variable action set is
published, and pretraining on some applications to test others has been
done on Android~\cite{zhao2024dinodroid}. What differs here is the oracle (verified seeded
faults instead of console errors), the outcome measure (time to the first verified fault
instead of coverage), and the question asked of the learner (does it transfer, and does it
use context).

\section{Language models in testing and as reward}
\label{sec:bg-llm}

\paragraph{LLMs as testers.} GPTDroid frames mobile GUI testing as a
question-answering loop with an LLM. DroidAgent sets its own task
goals and pursues them, reporting 61\% activity coverage against 51\% for earlier
techniques on 15 apps. AXNav executes natural-language accessibility
tests, and VETL uses a vision-language model for text inputs and element
selection on the web, with exploration guided by a multi-armed bandit. The 2026 study
tempers this line: its LLM-based web tester completed three to four times fewer actions in
the same time budget, because every decision invokes the model, and it repeated failed or
invalid actions~\cite{liu2026webgui}. Web-navigation benchmarks show how hard it is to act
reliably on unseen sites: in WebArena the best GPT-4 agent reached
14.41\% task success against 78.24\% for humans, and on Mind2Web
step success fell from 52.0\% on unseen tasks to 38.9\% on unseen websites.

\paragraph{LLMs as reward.} Kwon et al. prompt an LLM as a proxy
reward function. ELLM rewards an agent for achieving goals an LLM suggests,
motivated by the observation that novelty bonuses help little where most novelty is
irrelevant. Motif distils LLM preferences into an intrinsic
reward; Eureka and Text2Reward have the LLM
write reward code. Vision-language models have served as success
detectors and zero-shot reward models, whose
failures trace to the model's capability limits. Klissarov et al.
find that LLMs are better used to build reward models than to act directly.

\paragraph{Shaping.} Ng, Harada and Russell prove that adding a reward
equal to the discounted difference of a potential function over successive states,
$\gamma\Phi(s') - \Phi(s)$, leaves the optimal policy unchanged.
Wiewiora shows the equivalence to initialising Q-values from the
potential, and Devlin and Kudenko extend the guarantee to
potentials that change during learning. LLM-derived potentials already exist: Lin
et al. use score differences from an LLM-trained state scorer as a
potential, and report that with weaker models a direct reward outperformed the potential
form. M\"uller and Kudenko learn the potential from preferences given
by small vision-language models.

\paragraph{For this work.} Policy invariance concerns the optimal policy. It says nothing
about behaviour within a finite budget, which is what a tester is measured on. The evidence
for LLM rewards comes from games, robotics and synthetic web tasks; none of the retrieved
works evaluates an LLM-derived reward for testing. A small, local judge is on the weak side
of the capability evidence.

\section{Deep reinforcement learning and exploration}
\label{sec:bg-drl}

\paragraph{Value-based deep RL.} Q-learning learns the value
$Q(s,a)$ of taking action $a$ in state $s$ and acting well afterwards, by bootstrapping
from the next state's value. DQN~\cite{mnih2015human}
approximates $Q$ with a neural network. Double DQN~\cite{vanhasselt2016double} reduces
DQN's overestimation by decoupling action selection from evaluation. Dueling
networks separate state value from action advantage and give better
estimates when many actions have similar values, which fits pages with many near-equivalent
links. Prioritised replay replays informative transitions more
often, and Rainbow combines six such extensions, including
multi-step targets. These components were validated mainly on Atari, and Henderson
et al. show how variance makes deep RL improvements hard to
judge.

\paragraph{Variable action sets.} A web page offers a different set of actions on every
page. Dulac-Arnold et al. embed actions in a continuous space so
that a policy can generalise over them. The deep reinforcement relevance
network~\cite{he2016drrn} embeds state and text action separately and combines them to
approximate $Q$; this is the direct ancestor of the action-conditioned network
$Q(s,\phi(a))$ used here. Jain et al. treat generalisation to
unseen actions as its own problem, needing dedicated training objectives. Sentence
embeddings from compact distilled encoders
supply action-label features.

\paragraph{Exploration.} Count-based exploration rewards rarely visited states:
pseudo-counts generalise counting to large state spaces, and
hashing-based counts work with simple hash functions, provided
the hash has the right granularity. That is the same sensitivity to state abstraction that
the web-testing study measured~\cite{liu2026webgui}. Random network
distillation derives a novelty bonus from prediction error.
Go-Explore~\cite{ecoffet2021first} attributes poor exploration to two failures: forgetting
how to reach promising states, and not returning to a state before exploring from it. It
remedies both by keeping an archive of visited states and returning to them first.

\paragraph{Bandits.} A multi-armed bandit chooses repeatedly among options with unknown
payoff. Thompson sampling~\cite{russo2018tutorial} keeps a posterior
over each option's payoff, samples from it, and picks the best sample, which balances trying
uncertain options against exploiting good ones. With a Beta posterior over a success
probability, it needs only two counts per option, so it can learn within the few dozen
decisions of a single test run.

\paragraph{Generalisation.} Cobbe et al. find that agents
overfit to surprisingly large training sets, and Procgen shows that
diverse training distributions are essential. The survey by Kirk et al.
treats the number of training contexts as a core property of an evaluation protocol and
suggests fast online adaptation as a direction. With five training applications per fold,
the literature predicts weak transfer.

\section{Evaluating testing tools}
\label{sec:bg-eval}

Arcuri and Briand found that software-engineering studies of
randomised algorithms often did not account properly for randomness, and argue for enough
runs and rigorous statistical tests. Klees et al.~\cite{klees2018evaluating} examined 32
fuzzing papers and found problems in every evaluation, recommending multiple trials,
statistical tests, ground truth and a diverse benchmark; a follow-up over 150 papers found
the guidelines still widely ignored. In deep RL, Agarwal
et al. recommend stratified-bootstrap interval estimates over
point estimates.

When a tool may not find a bug within the time allowed, time to bug is a right-censored
outcome. Survival analysis handles this: the Kaplan--Meier
estimator estimates the probability of not yet having found
the bug at each time, and the restricted mean survival time (RMST)
averages the time to event up to a horizon. Magma applies Kaplan--Meier to time to
bug, and Herrera et al. report RMST with
confidence intervals and compare with the log-rank test. None of the GUI-testing papers
reviewed uses censoring-aware analysis: WebExplor averages 15 thirty-minute runs and applies
a Mann--Whitney U test~\cite{zheng2021webexplor}.

\section{Research gap}
\label{sec:bg-gap}

Within the literature reviewed:

\begin{enumerate}
  \item \textbf{Ground-truth evaluation of RL web testers.} Web RL testers are scored on
        coverage and console errors~\cite{zheng2021webexplor, gu2025webrled}. No retrieved work evaluates a web-testing policy, or
        its transfer to held-out applications, against a deterministic oracle for seeded
        faults.
  \item \textbf{An LLM-derived reward in testing.} LLM rewards and potentials exist for games
        and control, and LLMs in GUI testing act as
        the policy. Using an LLM
        only to shape a cheap learned policy is untested.
  \item \textbf{Censoring-aware, fault-centric evaluation in GUI testing.} Time to the first
        verified fault with right-censoring imports fuzzing's method into a field that
        reports averaged coverage, which is weakly related to failure
        discovery~\cite{liu2026webgui}.
  \item \textbf{Whether learning is what helps.} Humanoid's authors found their learned model
        was not the main cause of their gain~\cite{li2019humanoid}. Separating the
        contribution of learning from that of structure needs ablations against strong
        non-learning baselines.
\end{enumerate}

The contribution of this thesis is the combination and the evaluation discipline. The agent
architecture, the curiosity reward, the archive and the LLM potential each have direct
precedents and are not claimed as novel.
